\documentclass{article}

%%% Package List

\usepackage{datetime} % Dates.
\usepackage{fancyhdr} % Headers.
\usepackage{lastpage} % Page References.

\usepackage{graphicx} % Inserting Images.
\usepackage{float} % Postitioning Images.

\usepackage{dirtytalk} % Speech & Quote Marks.
\usepackage{hyperref} % Hyperlinks and References.
\usepackage{xcolor} % Coloured Text.
\usepackage{enumitem} % Letter Bullets. 

\usepackage{amssymb} % Maths Symbols.
\usepackage{amsmath} % Maths Tools.

\usepackage{algpseudocode} % Pseudocode.
\usepackage{algorithm} % Algorithms.
\usepackage{minted} % Code Snippets.

\usepackage[margin=1in]{geometry} % Reduces the Page Margins.

%%%



%%% Page Variables
\newcommand{\thetitle}{Computer Systems} % Insert the document title here.
\newcommand{\thesubtitle}{COMP1071: Durham University} % Insert the document subtitle here.
\newcommand{\theauthor}{Matthew Emerson} % Insert the name(s) of the author(s) here.
\newcommand{\thedate}{\monthyeardate\today} % Replace this if you want a static date.
%%%



%%% Custom Formats and Commands

\hypersetup{colorlinks=true, urlcolor=black, linkcolor=black}


% Date Formats.
\newdateformat{monthyeardate}
{%
    \monthname[\THEMONTH] \THEYEAR
}

\newdateformat{yeardate}
{%
    \THEYEAR
}


% Colouring and Spacing.
\newcommand{\colour}[2] {\color{#1}#2\color{black}}
\newcommand{\separate}[1] {\vspace{#1 pt}\noindent}


% Maths Number Sets.
\newcommand{\R}{\mathbb{R}}
\newcommand{\N}{\mathbb{N}}
\newcommand{\C}{\mathbb{C}}
\newcommand{\Q}{\mathbb{Q}}
\newcommand{\Z}{\mathbb{Z}}
\newcommand{\I}{\mathbb{I}}


% Custom Formats.
\newcommand{\definition}[2]{\textbf{Definition #1}: #2\label{def:#1}}
\newcommand{\example}[2]{\textbf{Example #1}: #2\label{ex:#1}}
\newcommand{\conjecture}[2]{\textbf{Conjecture #1}: #2\label{cnj:#1}}
\newcommand{\numberedentry}[3]{\textbf{#2 #1}: #3\label{entry:#1}}
\newcommand{\proof}[2]{\textit{Proof}: #2 \begin{flushright}\(\square\)\end{flushright}\label{prf:#1}}
\newcommand{\theorem}[3]{\textbf{Theorem #1}: #2 \par\separate{3}\proof{thm:#1}{#3}\label{thm:#1}}
\newcommand{\lemma}[3]{\textbf{Lemma #1}: #2 \par\separate{3}\proof{lem:#1}{#3}\label{lem:#1}}


% Advanced Formats.
\newcommand{\image}[3]
{\begin{figure}[H]
    \centering
    \includegraphics[width=#3\linewidth]{#1}
    \caption{#2}
    \label{fig:#2}
\end{figure}}

\newcommand{\algorithmBlock}[4]
{\begin{algorithm}
    \caption{#1}
    \begin{algorithmic}
        \State \textbf{Input}: #2
        \State \textbf{Output}: #3
        \separate{5}#4
    \end{algorithmic}
\end{algorithm}}

%%%



%%% Document Setup

% Title Page.
\title{\thetitle\\[0.5em]\large{\thesubtitle}}
\author{\theauthor}
\date{\thedate} 


% Custom Page Header.
\fancyhead{}
\fancyhead[R]{\thetitle\text{ - }\thesubtitle}

% Custom Page Footer.
\fancyfoot{}
\fancyfoot[L]{Written by \theauthor \\ 
\textit{\LaTeX \,Template: Copyright \date{\yeardate\today}, Matthew Emerson}} % Do not edit.
\fancyfoot[R]{Page \thepage \, of\, \pageref{LastPage}}


% Setting Up Pages.
\pagestyle{empty}\begin{document}
\maketitle\thispagestyle{empty} % Title Page.
\newpage\tableofcontents % Contents Page.

\clearpage\setcounter{page}{1}\newpage % Excludes the title and contents page from page count.
\pagestyle{fancy}

%%%



%%% Document Content

\part{Digital Electronics}\label{part:1.de}

\section{Transistors and Logic Gates}

\subsection{Transistors}

\definition{1.1.1}{\underline{Transistors} are electronically controlled switches that turn on or off when voltage or current is applied to a control terminal.}

\separate{5}\definition{1.1.2}{An \underline{nMOS transistor} is a sandwich of layers of conducting and insulating material. The gate is a conducting layer, on top of an insulating layer, on top of the substrate which is a semi-conductor.}

\image{}{}{}

\subsection{Logic Gates}

\newpage\part{Machine Architecture}\label{part:2.ma}



\newpage\part{Databases}\label{part:3.db}



\newpage\part{Operating Systems}\label{part:4.os}



\newpage\part{References}

\section{Module Professors}

The below is accurate as of the 2025/26 academic year and will not be updated in subsequent years.
\begin{itemize}
    \item \hyperref[part:1.de]{\textit{Digital Electronics}} and \hyperref[part:2.ma]{\textit{Machine Architecture}} are taught by Ioannis Ivrissimtzis.
    \begin{itemize}
        \item \textbf{Office Location}: Room 2023, MCS Building
        \item \textbf{Preferred Contact Methods}: ioannis.ivrissimtzis@durham.ac.uk
        \item \textbf{Pronouns}: he, him
    \end{itemize}
    
    \item \hyperref[part:3.db]{\textit{Databases}} is taught by Peter Davies-Peck.
    \begin{itemize}
        \item \textbf{Office Location}: Room 2063, MCS Building
        \item \textbf{Preferred Contact Methods}: peter.w.davies@durham.ac.uk
        \item \textbf{Pronouns}: he, him
    \end{itemize}
    
    \item \hyperref[part:4.os]{\textit{Operating Systems}} is taught by Anish Jindal.
    \begin{itemize}
        \item \textbf{Office Location}: Room 2065, MCS Building
        \item \textbf{Preferred Contact Methods}: anish.jindal@durham.ac.uk
        \item \textbf{Pronouns}: he, him
    \end{itemize}
\end{itemize}

\section{Cited Works}

\end{document}

%%%